% ============================================================
\section{Background and Related Work}
\label{sec:related-work}

We position M\textsuperscript{2} at the intersection of four literatures: TradFi capital-return mechanisms, DeFi tokenomic protocols, formal methods and mechanism design for automated market makers, and the economics of stablecoin pegs and runs. Each literature contributes a vocabulary and a set of negative results that the M\textsuperscript{2} design responds to.

\subsection{TradFi Capital-Return Mechanisms}
\label{sec:related-tradfi}

Three families of TradFi instruments inform the M\textsuperscript{2} design directly.

\paragraph{Buybacks and dividend policy.} Share repurchases have substituted for cash dividends as the dominant capital-return channel in U.S.\ equities since the late 1980s \citep{GrullonMichaely2002}. The substitution is partial: buybacks track earnings more flexibly than dividends \citep{Skinner2008}, but that flexibility is precisely what weakens their commitment value, and \citet{IkenberryLakonishokVermaelen1995} document a resulting market underreaction to repurchase announcements. M\textsuperscript{2} translates the buyback mechanic into an automated and deterministic form: a fixed 50\% slice of each revenue inflow is routed to an on-chain market buy whose execution is enforced by the protocol contract rather than by board discretion. There is no announcement-execution gap. Furthermore, by burning rather than retiring to treasury, the protocol guarantees that bought-back tokens cannot be re-issued---a guarantee that no public-equity buyback program can credibly make.

\paragraph{Closed-end funds and NAV redemption.} The persistent discount-to-NAV of closed-end funds (CEFs) is the canonical asset-pricing anomaly that maps most directly to M\textsuperscript{2} \citep{LeeShleiferThaler1991,Pontiff1996}. CEFs publish a net asset value but provide no redemption mechanism: investors who want out must sell on the secondary market, where the discount can persist for decades. \citet{CherkesSagiStanton2009} model the discount as a liquidity premium; \citet{BerkStanton2007} attribute part of it to management-fee capitalization. M\textsuperscript{2} is structurally a CEF whose NAV-redemption mechanism is enforced at the contract level: the discount cannot persist because any rational arbitrageur faced with $\Spot < \Floor$ buys on the AMM and redeems against the treasury. What M\textsuperscript{2} adds beyond NAV-redemption is \emph{monotone NAV}---in a CEF, the NAV itself fluctuates with portfolio mark-to-market, whereas in M\textsuperscript{2}, the contractually-defined $\Floor = \Tres/\Sup$ is non-decreasing by Theorem~\ref{thm:floor-monotonicity}.

\paragraph{Retractable bonds and other put options.} \citet{BrennanSchwartz1977} analyzed retractable bonds---fixed-income instruments that grant the holder a put option to redeem at par. The optionality is analytically equivalent to M\textsuperscript{2}'s \texttt{redeem} function: the holder may exit at a known price independent of secondary-market conditions. The differences are that retractable bonds redeem at a fixed schedule of par values, whereas M\textsuperscript{2} redeems at a state-dependent (but monotonically rising) floor, and that bond covenants can be modified by majority of holders whereas M\textsuperscript{2}'s redemption is immutable.

\paragraph{Mutual-fund redemption fragility.} \citet{ChenGoldsteinJiang2010} document strategic complementarities in open-end mutual fund redemptions: when illiquid-asset funds face redemption pressure, early redeemers extract value at the expense of remaining holders, generating a first-mover advantage that can produce run dynamics. M\textsuperscript{2}'s Theorem~\ref{thm:redemption-fairness} (Section~\ref{sec:formal-invariants}) is precisely the negation: redemption preserves the floor exactly, so there is no first-mover advantage in floor terms. The classical \citet{DiamondDybvig1983} bank-run equilibrium fails to obtain in M\textsuperscript{2} for the same structural reason.

\paragraph{Other reference instruments.} REIT mandatory distributions (the U.S.\ Internal Revenue Code \S856 requires REITs to distribute $\geq 90\%$ of taxable income to retain pass-through status; see \citet{Hartzell2014} for a study of institutional governance in this setting) provide a TradFi analogue to M\textsuperscript{2}'s hardcoded 50/50 revenue split, with the difference that REIT distributions can be retained at the cost of losing tax-advantaged status, whereas M\textsuperscript{2}'s split is a contract constant. Dividend reinvestment plans \citep{ScholesWolfson1989} compound returns into share accumulation, in contrast to M\textsuperscript{2}'s compounding into per-token floor. Money-market funds redeeming at par \citep{KacperczykSchnabl2013} provide the closest peg analogue, though M\textsuperscript{2} has no peg and instead has a monotone floor.

\paragraph{Synthesis.} Table~\ref{tab:tradfi-analogues} summarizes the mapping. No single TradFi instrument composes the three features that M\textsuperscript{2} composes (the headline framing in Section~\ref{sec:introduction} names them); the composition is enabled by the contract-as-trust substrate rather than by any individual feature being novel.

\begin{table}[t]
\centering
\small
\caption{TradFi capital-return instruments and their M\textsuperscript{2} analogues.}
\label{tab:tradfi-analogues}
\setlength{\tabcolsep}{4pt}
\begin{tabular}{p{2.6cm} p{4.2cm} p{4.4cm} p{3.4cm}}
\toprule
\textbf{TradFi instrument} & \textbf{Key feature} & \textbf{M\textsuperscript{2} analogue} & \textbf{What M\textsuperscript{2} adds} \\
\midrule
Closed-end fund & NAV is published but not enforced; persistent discount & Floor $\Floor = \Tres/\Sup$ enforced by \texttt{redeem} & NAV itself is monotone-nondecreasing \\
Open-market buyback & Sponsor-discretionary; announcement-execution gap & Hardcoded 50\% revenue split to buy-and-burn & No discretion; bought-back tokens are burned, not retired \\
Retractable bond & Holder put at par on schedule & \texttt{redeem} at floor, on demand & State-dependent rising floor; immutable terms \\
DRIP & Compounding via share accumulation & Compounding via floor growth & Holder needs no action; compounding is automatic \\
REIT distribution & Required by tax structure & Required by contract constant & No legal entity needed \\
Tobin tax / Sec.\ 31 fee & Externality charge to state & 3\% sell fee accruing to floor & Pigovian charge recycled to common pool \\
\bottomrule
\end{tabular}
\end{table}

\subsection{DeFi Tokenomic Protocols}
\label{sec:related-defi}

\citet{Werner2022SoK} provide a taxonomic survey of decentralized finance that frames the landscape we situate M\textsuperscript{2} within. We comment specifically on prior protocols that have engaged with redemption, treasury backing, or perpetual buybacks.

\paragraph{OlympusDAO and forks.} OHM \citep{OlympusDAO2021} introduced the ``backed not pegged'' framing and the concept of risk-free-value (RFV) per token. The framing was rhetorical rather than enforced: holders could not redeem OHM against the RFV treasury, and treasury composition could be changed by governance vote. The empirical record of OHM forks (Wonderland, Klima, and others) is dominated by treasury raids and confidence collapses. M\textsuperscript{2} responds to this lineage by (i) making redemption a first-class contract function, (ii) removing governance entirely, and (iii) holding only a single transparent stablecoin in treasury. Crucially, M\textsuperscript{2}'s supply only decreases over time, in contrast to OHM's staking-rebase inflation.

\paragraph{Reflexer Finance / RAI.} \citet{Reflexer2021} introduced a PID-controlled redemption price that can diverge from market price. RAI and M\textsuperscript{2} share the conceptual move of separating ``redemption price'' from ``market price,'' but RAI's redemption price is mean-reverting around a control target whereas M\textsuperscript{2}'s is monotone-nondecreasing.

\paragraph{Frax, Liquity, MakerDAO.} \citet{Frax2021} introduced fractional-algorithmic stablecoins. \citet{Liquity2021} demonstrated a governance-free, immutable, ETH-collateralized stablecoin (LUSD) and is the closest prior art for M\textsuperscript{2}'s ``immutable plus no governance'' design stance. \citet{Maker2019} pioneered multi-collateral, governance-rich CDP-style redemption. Each of these systems is a stablecoin (pegged or controlled); M\textsuperscript{2} is not a stablecoin. The redemption mechanism is shared in structure, but M\textsuperscript{2}'s redemption is to a state-dependent floor rather than to a target peg.

\paragraph{Bonding curves and protocol-owned liquidity.} \citet{Bancor2017} introduced bonding curves as a primitive for protocol-enforced redemption pricing; M\textsuperscript{2}'s $\Floor = \Tres/\Sup$ is a degenerate bonding curve where the per-token payout is independent of position on the curve. \citet{Heimbach2022} characterize liquidity-provider entry and exit behavior in automated market makers, a structural sibling of the protocol-owned-liquidity regime M\textsuperscript{2} adopts; M\textsuperscript{2} inherits this lineage but locks the position permanently.

\subsection{Formal Methods and Mechanism Design for AMMs}
\label{sec:related-formal}

\citet{Angeris2020} provide the canonical formal treatment of constant-function market-maker properties; \citet{Angeris2021} extend this to no-arbitrage results for Uniswap. \citet{Bartoletti2021} develop a theory of automated market makers that informs the proof style we adopt in Section~\ref{sec:formal-invariants}. \citet{UniswapV4} specify the v4 hook architecture that M\textsuperscript{2} uses for asymmetric per-swap fee overrides via \texttt{beforeSwap} (with \texttt{LPFeeLibrary.OVERRIDE\_FEE\_FLAG} on dynamic-fee pools).

On the adversarial side, \citet{Daian2020} introduced the modern study of maximum extractable value, and \citet{Qin2022} provide empirical magnitudes. \citet{HeimbachWattenhofer2022} apply game-theoretic mechanism design to sandwich attack prevention; the principles inform M\textsuperscript{2}'s anti-MEV mitigations (private orderflow, randomized delay, TWAP execution). \citet{Babel2023} introduce a methodology for adversary-bounded value extraction proofs in smart-contract economic-security analysis, which is the proof genre our Section~\ref{sec:economic-analysis} bank-run analysis lives in.

A distinct strand of recent work concerns \emph{Loss-Versus-Rebalancing} (LVR) for AMM liquidity providers \citep{MilionisMoallemiRoughgarden2023}. The standard LVR framing presents the LP-versus-external-CEX price differential as a cost extracted by arbitrageurs. M\textsuperscript{2} inverts this framing structurally: the LP position is owned by the protocol, the protocol's value flows into the floor, and---under the asymmetric fee mechanism---LVR-equivalent value is plausibly captured rather than lost. We conjecture this inversion as Conjecture~\ref{conj:lvr-capture} and discuss the unit-reconciliation gap that prevents promoting it to a theorem in Section~\ref{sec:lim-open-math}.

For broader context on formal verification, \citet{Tolmach2022} survey smart-contract specification and verification methods. The skeptical view of decentralized market-making is articulated in \citet{Park2023}; we engage with this critique in Section~\ref{sec:limitations} given that M\textsuperscript{2}'s LP is internal rather than external.

\subsection{Stablecoin Run Dynamics}
\label{sec:related-stablecoins}

\citet{KlagesMundt2020} provide an economic taxonomy of stablecoin designs and risk-based models. The Terra--Luna collapse \citep{LiuMakarovSchoar2023,Briola2023} is the benchmark case study for how endogenous-collateral algorithmic stablecoins enter a reflexive death spiral under coordinated redemption pressure. We discuss M\textsuperscript{2}'s structural disanalogy with Terra in Section~\ref{sec:economic-analysis}: M\textsuperscript{2}'s collateral is exogenous (a deployer-chosen USD-pegged stablecoin) and one-way-in (the only path stable leaves the treasury is \texttt{redeem}, which is in lockstep with supply burn). The Terra-style spiral is structurally absent.

Adjacent work on peg dynamics includes \citet{Gorton2022} on leverage and stablecoin pegs, \citet{Catalini2022} on the economic foundations of stablecoins, \citet{Routledge2022} on currency stability via blockchain, and \citet{dAvernas2023} on the broader question of stablecoin stability under reserve composition. \citet{Eichengreen2019} situates the debate historically. M\textsuperscript{2} is explicitly not a stablecoin---it has no peg target, no e-money classification under MiCA's tests, and a deliberately variable spot price. Nonetheless the analytical tools of the stablecoin literature transfer because the redemption mechanic poses the same liquidity questions as a depeg event.
